% Created with jtex v.1.0.20
\documentclass{article}
\usepackage{hyperref}
\usepackage{datetime}
\usepackage{graphicx}
\usepackage{natbib}
\bibliographystyle{abbrvnat}

%%%%%%%%%%%%%%%%%%%%%%%%%%%%%%%%%%%%%%%%%%%%%%%%%%
%%%%%%%%%%%%%%%%%%%%  imports  %%%%%%%%%%%%%%%%%%%
\usepackage{amsmath}
%%%%%%%%%%%%%%%%%%%%%%%%%%%%%%%%%%%%%%%%%%%%%%%%%%


% colors for hyperlinks
\hypersetup{colorlinks=true, allcolors=blue}

\newcommand{\logo}{
  \href{https://curvenote.com}{\includegraphics[width=2cm]{curvenote.png}}
}

\title{Lecture 3 Notes - Characteristics of Time Series}

\author{Liberty Hamilton}

\newdate{articleDate}{27}{1}{2026}
\date{\displaydate{articleDate}}

\begin{document}
\maketitle
\begin{center}\logo\end{center}


\begin{itemize}
\item \textbf{Reading}: Chapter 1.3-1.7 -- Shumway and Stoffer
\end{itemize}

\section{Review / basic concepts}

\begin{itemize}
\item Stochastic process
\item White noise
\item Moving average
\item Autoregression
\item Random walk (w/ and w/o drift)
\item Signals in noise
\end{itemize}

\section{This week - Measures of dependence}

\begin{itemize}
\item Mean function
\item Autocovariance
\item Autocorrelation
\item Cross-covariance
\item Cross-correlation
\item Stationarity
\item Multidimensional time series
\item Random number generation
\end{itemize}

\section{Mean and variance}

Last week and in your lab, we saw the concept of how mean and variance differ for a white noise process, a random walk, and a random walk with drift. The \textit{mean} and \textit{variance functions} of a time series are useful descriptors because they helps us determine something about the drift and the spread of the data that we should expect over time.

The \textit{mean function} is defined as $\mu_{xt} = \mathbb{E}(x_t)$

The \textit{variance function} is $\sigma^2_t = \operatorname{Var}(x_t) = \mathbb{E}[(x_t -\mu_t)^2]$.

Let's revisit some examples from before:

\subsection{White noise}

For a white noise time series, $\mu_{wt} = \mathbb{E}(w_t)=0$ for all $t$. $\operatorname{Var}(w_t) = 1$ (for a Gaussian white noise series).

\subsection{Moving average}

What if we apply a 3-point moving average? Although this induces some correlation structure, it actually does not change the mean function at all:

$\mu_{vt} = \mathbb{E}(v_t) = \frac{1}{3}[\mathbb{E}(w_{t -1})+\mathbb{E}(w_{t})+\mathbb{E}(w_{t+1})] = 0$

\subsection{Random walk with drift}

For the random walk with drift, the mean function is just the line $\mu_{xt} = \delta t + \displaystyle\sum_{j=1}^t E(w_j) = \delta t$.

\subsection{Signal plus noise}

And for a signal plus noise (as an example, we'll use this sinusoid):

\begin{equation}
\begin{aligned}
\mu_{xt} = \mathbb{E}(x_t) &= \mathbb{E} [A \cos(2\pi \omega t + \phi) + w_t] \\ 
&= \mathbb{E} [A \cos(2\pi \omega t + \phi)] + \mathbb{E}[w_t] \\
&= \mathbb{E} [A \cos(2\pi \omega t + \phi)]
\end{aligned}
\end{equation}

\subsection{Autocovariance}

What if we instead want to know something about the dependence between two points $s$ and $t$ within the same time series. We'll call this \textit{autocovariance}, $\gamma$.

$\gamma_x(s,t) = \operatorname{cov}(x_s,x_t) = \mathbb{E}[(x_s -\mu_s)(x_t -\mu_t)]$

When $s=t$, we have:

$\gamma_x(t,t) = \mathbb{E}[(x_t -\mu_t)^2] = \operatorname{Var}(x_t)$

For white noise, there should be no dependence between differing time points. By definition, $\mathbb{E}(w_t)=0$ and:

\begin{equation}
\gamma_w(s,t) = \operatorname{cov}(w_s,w_t) = 
\begin{cases}
\sigma^2_w, & \text{if } s=t, \\
0, & s \neq t
\end{cases}
\end{equation}

\subsection{Covariance of linear combinations}

If we have random variables $U=\displaystyle\sum_{j=1}^m a_j X_j$ and $V=\displaystyle\sum_{k=1}^r b_k Y_k$ that are linear combinations of (finite variance) random variables ${X_j}$ and ${Y_k}$, then the covariance of these is:

$\operatorname{cov}(U,V) = \displaystyle\sum_{j=1}^m \displaystyle\sum_{k=1}^r a_j b_k \operatorname{cov}(X_j, Y_k)$

Also, $\operatorname{var}(U) = \operatorname{cov}(U,U)$.

So how can we use this? Now we can try this for the moving average example.

\begin{equation}
\begin{aligned}
\gamma_v(s,t) = \operatorname{cov}(v_s, v_t) &= \operatorname{cov}(\tfrac{1}{3}(w_{s-1}+w_s+w_{s+1}), \tfrac{1}{3}(w_{t-1}+w_t+w_{t+1})) \\
&= \tfrac{1}{9} \operatorname{cov}(w_{s-1}+w_s+w_{s+1}, w_{t-1}+w_t+w_{t+1})
\end{aligned}
\end{equation}

When $s=t$, we have:
\begin{equation}
\begin{aligned}
\gamma_v(t,t) &= \operatorname{cov}(v_t, v_t) \\
&= \operatorname{cov}(\tfrac{1}{3}(w_{t-1}+w_t+w_{t+1}), \tfrac{1}{3}(w_{t-1}+w_t+w_{t+1})) \\
&= \tfrac{1}{9} (\operatorname{cov}(w_{t-1},w_{t-1}) + \operatorname{cov}(w_{t},w_{t}) + \operatorname{cov}(w_{t+1},w_{t+1}) \\
&\quad + 2\operatorname{cov}(w_{t-1}, w_t)) + 2\operatorname{cov}(w_{t}, w_{t+1})) + 2\operatorname{cov}(w_{t-1}, w_{t+1}))\\
&= \tfrac{1}{9} (\sigma_w^2 + \sigma_w^2 + \sigma_w^2 + 0 +0 +0)\\
&= \tfrac{3}{9} \sigma_w^2 \\
&= \tfrac{1}{3} \sigma_w^2
\end{aligned}
\end{equation}

When $s=t+1$, we have:
\begin{equation}
\begin{aligned}
\gamma_v(t+1,t) &= \operatorname{cov}(\tfrac{1}{3}(w_{t}+w_{t+1}+w_{t+2}), \tfrac{1}{3}(w_{t-1}+w_t+w_{t+1})) \\
&= \tfrac{1}{9}(\operatorname{cov}(w_t,w_t)+\operatorname{cov}(w_{t+1},w_{t+1}))\\
&= \tfrac{2}{9}\sigma^2_w
\end{aligned}
\end{equation}

If we then follow this for more lags, we get:

\begin{equation}
\gamma_v(s,t) = 
\begin{cases}
\tfrac{3}{9}\sigma^2_w, & \text{if } s=t, \\
\tfrac{2}{9}\sigma^2_w, & \text{if } |s-t|=1, \\
\tfrac{1}{9}\sigma^2_w, & \text{if } |s-t|=2, \\
0, & \text{if } |s-t|>2
\end{cases}
\end{equation}

Why is this interesting? The autocovariance depends only on the \textit{lag} between $s$ and $t$ and not on the absolute location of these time points. We'll come back to this when we talk about \textit{stationarity}.

\subsection{Autocovariance of a random walk}

Recall that a random walk (with or without drift) is:

$x_t = \delta t + \displaystyle\sum_{i=1}^t w_i$

So,
\begin{equation}
\begin{aligned}
\gamma(s,t) &= \operatorname{cov}(x_s, x_t) \\
&= \operatorname{cov}(\delta s + \displaystyle\sum_{i=1}^s w_i, \delta t + \displaystyle\sum_{i=1}^t w_i)\\
&= \sigma^2 \min(s,t)
\end{aligned}
\end{equation}

In this case, the autocovariance \textit{does} depend on the particular $s$ and $t$ chosen.

What if we want a bounded measure?

\subsection{Autocorrelation function}

We can calculate the autocorrelation function (ACF) as:

$\rho(s,t) = \frac{\gamma(s,t)}{\sqrt{\gamma(s,s)\gamma(t,t)}}$

This measures the linear predictability of the time series at time $t$ ($x_t$) using $x_s$. We can also extend this to looking at the linear predictability of one time series to another, by extending into the concepts of \textit{cross-covariance} and \textit{cross-correlation}.

\subsection{Cross-covariance}

Between two time series $x_t$ and $y_t$:

$\gamma_xy(s,t) = \operatorname{cov}(x_s,y_t) = \mathbb{E}[(x_s -\mu_{xs})(y_t -\mu_{yt})]$

This tells us how the values in $y$ relate to the values in $x$ over time.

Let's think about a simple example:

$y_t = x_{t -2}$

What is $\gamma_{xy}(k)$ (for lag $k$)? At what lag is $\gamma_{xy}$ maximized?

We can also have the normalized version:

\subsection{Cross-correlation}

$\rho_{xy}(s,t) = \frac{\gamma_xy(s,t)}{\sqrt{\gamma_x(s,s)\gamma_y(t,t)}}$

This is bounded such that $-1 \leq \rho(s,t) \leq 1$

\textit{Poll example}

\section{Next time:}

\begin{itemize}
\item Stationarity
\item Calculating sample correlation/covariance
\item Multi-D time series
\item RNGs
\item Measures of dependence! Read SS Chapter 1, sections 1.3-1.7.
\end{itemize}


\end{document}
